\section{优先研究方向与实验设计}
\label{sec:future-directions}

以下方案由既有证据和模型候选共同提出，均为实验假设，不是已经验证的目标相配方。建议按“局部相图--占位/氧计量--晶体生长--类比替代”的顺序推进；前一阶段的相组成结果决定后一阶段是否值得开展。

\subsection{方向一：目标组成附近的局部相图}

第一优先级是在$\mathrm{Ba_5Y_{12}ZnSi_8O_{40}}$周围建立小步长组成网格。基础前驱体采用模型排名第一且化学上常规的BaCO$_3$、Y$_2$O$_3$、ZnO和SiO$_2$；目标化学计量对应5:6:1:8的摩尔比。本方向属于模型预测，待实验验证。建议先按无Zn固相研究的温区分段筛查：约1000\,$^\circ$C、12\,h进行初始反应，复磨后在1300\,$^\circ$C、24\,h退火，并依据PXRD决定是否重复；只有仍有明显未反应物且失重和坩埚相容性可控时，才增设1450--1600\,$^\circ$C的短时节点。温度和时间来自相近Ba--Y固相工作的筛查区间\cite{yamane2024microstructure,gulay2024navigation}，用于安排实验矩阵，不表示目标相必在这些条件下形成。

组成轴至少应包含(i)目标点附近的Ba/Y小幅偏离；(ii)Zn含量；(iii)与氧计量相关的无Zn参照相。每个样品均记录烧前和烧后质量、热史和冷却方式，并以定量PXRD/Rietveld、SEM--EDS/EPMA和必要的ICP分析给出相组成与实际元素比。若出现玻璃或多相边界，应在该区域加密采样，而不是只保留“最好看”的衍射谱。

\subsection{方向二：Zn--Y--氧计量耦合}

一个可检验的电荷平衡系列是
\[
\mathrm{Ba_5Y_{13-x}Zn_x[SiO_4]_8O_{8.5-x/2}},
\qquad 0\leq x\leq1.25,
\]
其中$x=0$对应已报道的无Zn四方参照相，$x=1$在元素计量上对应目标组成。建议以$x=0,0.25,0.50,0.75,1.00,1.25$为首轮离散点。本方向属于模型预测，待实验验证。前驱体仍采用BaCO$_3$、Y$_2$O$_3$和SiO$_2$，以ZnO调节Zn含量；该表达式只是电荷平衡假设，不预设Zn的真实占位或氧空位分布。判断应同时考察晶胞参数连续性、相分数、实际Zn/Y比和占位精修；若条件允许，中子衍射或独立的氧含量和价态测量可进一步区分Zn替位、氧缺陷与两相混合。

这一系列的科学价值高于单点复现：结构参数连续变化支持有限固溶，晶胞基本不变而两相比例变化更接近两相区，只在$x\approx1$出现新衍射峰则提示线化合物。三种结果都能约束目标相的形成机制。

\subsection{方向三：固相成相与高温溶液长晶并行}

固相支线用于定位批量相区，高温溶液支线用于获得可供SCXRD分析的晶体。本方向属于模型预测，待实验验证。高温溶液支线可把Ba--Y--Si--O中已验证的K$_2$CO$_3$/MoO$_3$助熔体系作为对照起点，前驱体采用方向一的BaCO$_3$、Y$_2$O$_3$、ZnO和SiO$_2$；在1100--1200\,$^\circ$C范围设置最高温度，保温约3\,h，以1--2\,K\,h$^{-1}$慢冷至约900\,$^\circ$C，并系统改变营养物与助熔剂比例以及Y$_2$O$_3$/SiO$_2$比\cite{kolitsch2009crystal,wierzbickawieczorek2017high}。这些数值来自无Zn同谱系晶体生长，必须与目标投料的固相支线并行比较，不能称为目标相处方。

每个长晶批次至少报告液相组成、装填率、最高温度、保温、冷速、分离方法、晶体尺寸和产率；晶体用SCXRD和EDS验证，剩余母体用PXRD和成分分析检查。若只得到少量目标单晶而母体远离目标组成，应解释为液相选择性生长，而不是批量相区。

\subsection{方向四：Mg/Co类比与模型前驱体}

本文调用本地两步无机逆合成模型，对三个元素式进行了Top-5组合预测。模型对Zn、Mg和Co目标的首选均保留BaCO$_3$、Y$_2$O$_3$和SiO$_2$，第四种原料分别为ZnO、MgO和Co$_3$O$_4$；三条首选在各自候选池内的概率分别为0.594、0.522和0.298。本方向属于模型预测，待实验验证。表~\ref{tab:mcp-experiments}把模型输出转成可检验的最小实验。模型只依据元素计量提出前驱体组合，不用于判断硅酸根结构单元；正文始终保留目标文献使用的结构式。

\begin{table}[t]
  \centering
  \caption{由两步无机逆合成模型提出的首选前驱体及实验用途。所有组合均为模型排序（未经过化学路线验证）；工艺条件来自相近体系文献，需经相图实验复核。}
  \label{tab:mcp-experiments}
  \footnotesize
  \setlength{\tabcolsep}{3.0pt}
  \begin{tabular}{P{1.55cm} P{3.6cm} P{1.2cm} P{6.7cm}}
    \toprule
    假设目标 & 模型首选前驱体 & 候选池概率 & 建议实验与判据\\
    \midrule
    Zn目标 & BaCO$_3$ + Y$_2$O$_3$ + SiO$_2$ + ZnO & 0.594 & 作为局部相图的基础投料；先固相筛查，再把近单相坐标转入助熔长晶；PXRD/Rietveld与SCXRD分别回答批量相区和结构身份。 \\
    Mg类比 & BaCO$_3$ + Y$_2$O$_3$ + SiO$_2$ + MgO & 0.522 & 以Mg$^{2+}$作为同价尺寸对照；沿与Zn相同的组成/热史矩阵筛查。只有出现新相且实际组成连续变化时，才讨论同一位点替代。 \\
    Co类比 & BaCO$_3$ + Y$_2$O$_3$ + SiO$_2$ + Co$_3$O$_4$ & 0.298 & 以空气为首轮条件，测定实际组成和相分数，避免将副相误判为结构替代。 \\
    \bottomrule
  \end{tabular}
\end{table}

模型给出的其余Zn路线包含额外BaO、BaF$_2$、Ba(NO$_3$)$_2$或非常规数据库标签，概率均明显低于首选，不应并列推荐。额外Ba源会改变局部相图，BaF$_2$增加含氟挥发物和废物流风险，非常规标签还可能反映训练数据标准化问题。Mg路线中的MgCO$_3$可作为单独的前驱体形态对照，但不应与MgO同时加入；Co的多种氧化物前驱体不应视为三个等价方案。

\subsection{能够促进科学发现的实验组织}

每轮实验应以“最大化可区分结论”为标准，而不是只追求得到晶体。最有信息量的组合是：(i)小步长组成网格配合定量相分析，建立局部相图；(ii)同一组成的固相和助熔平行实验，区分平衡相区与液相选择性成核；(iii)Zn含量与氧计量联动系列，检验位点和缺陷假设。所有负结果同样保存，因为它们决定下一轮组成选择并约束“可合成”的边界。
